\section*{Capítulo \ref{ch:b2:crustal-melting} --- A fusão da crosta: migmatitos e granitos}

\begin{solution}{exo:b2:crustal-melting:1}
(a) Por \qty{100}{\gram}: $n(\ce{Al2O3}) = 17.03/101.96 = \qty{0.16703}{mol}$,
$n(\ce{CaO}) = 0.09183$, $n(\ce{Na2O}) = 0.06783$, $n(\ce{K2O}) = 0.03076$.
A/CNK $= 0.16703/0.19042 = 0.877$; A/NK $= 0.16703/0.09859 = 1.694$.
(b) A/CNK $< 1 <$ A/NK: \omterm{def:b2:crustal-melting:asi}{metaluminoso}. O cálcio que os feldspatos não podem receber vai para
silicatos de cálcio: hornblenda (ou clinopiroxênio, titanita).
\end{solution}

\begin{solution}{exo:b2:crustal-melting:2}
(a) A muscovita e a granada contêm alumínio além dos feldspatos: \omterm{def:b2:crustal-melting:asi}{peraluminoso}; esse granito
pode ser um tipo S. (b) A hornblenda e a titanita contêm cálcio além dos feldspatos:
\omterm{def:b2:crustal-melting:asi}{metaluminoso}. (c) A egirina e um anfibólio sódico contêm sódio sem o alumínio
do feldspato: \omterm{def:b2:crustal-melting:asi}{peralcalino}.
\end{solution}

\begin{solution}{exo:b2:crustal-melting:3}
Um quilômetro de rocha pesa $2800\times9.81\times1000 = \qty{2.747e7}{Pa}$, logo
$P = 0.2747\,z$ (\unit{kbar}, $z$ em \unit{km}). A reta do solidus é $T = 680 - 7.5(P -
2) = 695 - 2.060\,z$. Igualando $10 + 20z = 695 - 2.060z$, obtém-se $z = 685/22.060 =
\qty{31.1}{km}$, $T = \qty{631}{\celsius}$, $P = \qty{8.5}{kbar}$.
\end{solution}

\begin{solution}{exo:b2:crustal-melting:4}
(a) As camadas brancas são \omterm{def:b2:crustal-melting:migmatite}{leucossomas}, as faixas negras \omterm{def:b2:crustal-melting:migmatite}{melanossomas}, o gnaisse comum o
mesossoma. (b) \omterm{def:b2:crustal-melting:migmatite}{Leucossomas}, \omterm{def:b2:crustal-melting:migmatite}{melanossomas} e mesossomas distintos e bem definidos: um \omterm{def:b2:crustal-melting:migmatite}{metatexito},
de estrutura estromática. (c) A rocha se tornou um \omterm{def:b2:crustal-melting:migmatite}{diatexito} (schlieren e nebulítico):
a fusão foi mais longe, os minerais escuros fundiram também, e o fundido permaneceu no lugar,
apagando o bandamento.
\end{solution}

\begin{solution}{exo:b2:crustal-melting:5}
(a) $M = 39.10 + 3\times24.31 + 26.98 + 3\times28.09 + 12\times16.00 + 2\times1.008 =
\qty{417.30}{\gram\per\mole}$; $w = 18.02/417.30 = \qty{4.32}{\percent}$.
(b) $F = 0.25\times0.0432/0.040 = 0.270$, ou \qty{27}{\percent} de fundido em massa: dentro dos
20 a \qty{30}{\percent} em volume produzidos pelos experimentos de \omterm{def:b2:crustal-melting:reactions}{fusão por desidratação} da biotita em
rochas semelhantes.
\end{solution}

\begin{solution}{exo:b2:crustal-melting:6}
(a) Feldspato potássico \qty{278.35}{\gram\per\mole}, albita 262.24, anortita 278.22, muscovita
398.33. (b) Quantidades: feldspato potássico $30/278.35 = \qty{0.10778}{mol}$, albita $0.13347$,
anortita $0.01797$, muscovita $0.01255$. O feldspato potássico é $\tfrac12$\ce{K2O.}$\tfrac12$\ce{Al2O3},
a albita $\tfrac12$\ce{Na2O.}$\tfrac12$\ce{Al2O3}, a anortita \ce{CaO.Al2O3}, a muscovita
$\tfrac12$\ce{K2O.}$\tfrac32$\ce{Al2O3} por unidade de fórmula. $n(\ce{Al2O3}) = 0.05389 +
0.06673 + 0.01797 + 0.01883 = \qty{0.15742}{mol}$; $n(\ce{K2O}) = 0.05389 + 0.00628 =
0.06017$, $n(\ce{Na2O}) = 0.06673$, $n(\ce{CaO}) = 0.01797$, soma $0.14487$.
A/CNK $= 0.15742/0.14487 = 1.087$. (c) Sem a muscovita, os três feldspatos dão
exatamente 1 (\cref{prop:b2:crustal-melting:feldspar}); a muscovita, com um \ce{Al2O3}
a mais que o feldspato potássico por unidade de fórmula, torna a rocha \omterm{def:b2:crustal-melting:asi}{peraluminosa}.
\end{solution}

\begin{solution}{exo:b2:crustal-melting:7}
Basalto: $900 + 2.5\times(1200 - 850) = \qty{1775}{\mega\joule\per\cubic\metre}$. Crosta:
$2.5\times(850 - 650) + 0.20\times900 = \qty{680}{\mega\joule\per\cubic\metre}$.
$V_c/V_b = 1775/680 = 2.6$. A estimativa ignora o calor que escapa por condução
para cima e para os lados e o tempo que o calor leva para se espalhar; ela supõe também que todo o
calor do basalto chega a uma crosta pronta para fundir: é um limite superior.
\end{solution}

\begin{solution}{exo:b2:crustal-melting:8}
(a) $n(\ce{Al2O3}) = 14.15/101.96 = 0.13878$, $n(\ce{CaO}) = 0.01753$, $n(\ce{Na2O}) =
0.07809$, $n(\ce{K2O}) = 0.04416$; A/CNK $= 0.13878/0.13978 = 0.993$, A/NK $=
0.13878/0.12225 = 1.135$: \omterm{def:b2:crustal-melting:asi}{metaluminoso}, por pouco. (b) Com \qty{14.30}{\percent} de \ce{Al2O3},
$0.14025/0.13978 = 1.003$; com \qty{14.00}{\percent}, $0.13731/0.13978 = 0.982$. O índice
fica de um lado e de outro de 1 dentro da incerteza certificada: a obsidiana fica no
limite metaluminoso--peraluminoso, e sua classe não é robusta.
\end{solution}

\begin{solution}{exo:b2:crustal-melting:9}
(a) Médias $\bar P = \qty{5.0}{kbar}$, $\bar w = \qty{10.045}{\percent}$;
$\sum(P - \bar P)(w - \bar w) = 17.62$, $\sum(P - \bar P)^2 = 14.0$; inclinação
$17.62/14.0 = \qty{1.259}{\percent}$ por quilobar, ordenada na origem $10.045 - 1.259\times5.0 =
\qty{3.75}{\percent}$. A inclinação concorda com os 1.27 indicados a menos de
\qty{1}{\percent}. (b) $3.75 + 1.259\times6 = \qty{11.30}{\percent}$. (c) $10/0.1130 =
\qty{88.5}{\gram}$ de fundido. (d) A reta se apoia em dados entre 3 e \qty{8}{kbar} a
\qty{800}{\celsius} apenas; nada garante que a solubilidade permaneça linear a pressões
mais altas, onde o fundido e o fluido se tornam mais parecidos.
\end{solution}

\begin{solution}{exo:b2:crustal-melting:10}
(a) A reta sobe \qty{4}{kbar} em \qty{50}{\celsius}: $P = 6 + 0.08\,(T - 750)$;
a \qty{770}{\celsius}, $P = \qty{7.6}{kbar}$. (b) De 7.6 a \qty{6}{kbar}, a
temperatura cai ao longo da reta de 770 a \qty{750}{\celsius}; o calor liberado,
$2.5\times20 = \qty{50}{\mega\joule\per\cubic\metre}$, funde $50/900 = 0.056$, cerca de
\qty{6}{\percent} da rocha em volume. (c) É a janela proposta para os
\omterm{def:b2:crustal-melting:leucogranite}{leucogranitos} do Himalaia, 6 a \qty{8}{kbar} e 750 a \qty{770}{\celsius}: a descompressão
sozinha produz fundido, mas apenas alguns por cento a cada travessia da curva.
\end{solution}

\begin{solution}{exo:b2:crustal-melting:11}
$\ell^2/\kappa = (2000)^2/10^{-6} = \qty{4e12}{s} = \qty{0.13}{Myr}$ para a \omterm{def:b2:magma-emplacement:bodies}{soleira};
$(10^4)^2/10^{-6} = \qty{1e14}{s} = \qty{3.2}{Myr}$ para a camada sobreplacada. Os dois são curtos
diante de dez milhões de anos: um único evento não pode manter a crosta fundida por tanto tempo, o que exige
calor fornecido repetidas vezes, por intrusões repetidas ou pelo aquecimento lento de uma
crosta espessada.
\end{solution}

\begin{solution}{exo:b2:crustal-melting:12}
(a) $5.0x + 2.0(1 - x) = 2.75$ dá $x = 0.75/3.0 = 0.25$: um quarto do neossoma é
\omterm{def:b2:crustal-melting:migmatite}{leucossoma}. (b) $0.25\times0.20 + 0.75\times4.0 = 0.05 + 3.00 = \qty{3.05}{\percent}$, como
no mesossoma: a interpretação em sistema fechado se sustenta. (c) Para \qty{100}{\gram} de
rocha original, formaram-se \qty{25}{\gram} de fundido e \qty{75}{\gram} de resíduo; se $m$ gramas de
fundido saíram, $(25 - m)/(100 - m) = 0.10$, logo $25 - m = 10 - 0.1m$ e $m = 15/0.9 =
\qty{16.7}{\gram}$: dois terços do fundido deixaram o afloramento.
\end{solution}

\begin{solution}{pb:b2:crustal-melting:1}
\noindent\textbf{Parte I.}

\textbf{1.} $39.10 + 3\times26.98 + 3\times28.09 + 12\times16.00 + 2\times1.008 =
\qty{398.33}{\gram\per\mole}$.

\textbf{2.} $18.02/398.33 = \qty{4.52}{\percent}$.

\textbf{3.} $200\times0.0452 = \qty{9.04}{\gram}$ de água por quilograma de xisto.

\textbf{4.} $9.97 + (13.82 - 9.97)\times2/3 = \qty{12.54}{\percent}$. O fundido mínimo é
$9.04/0.1254 = \qty{72.1}{\gram}$ por quilograma: $F \ge \qty{7.21}{\percent}$
(\cref{prop:b2:crustal-melting:water}).

\textbf{5.} $9.04/0.060 = \qty{151}{\gram}$: $F = \qty{15.1}{\percent}$.

\textbf{6.} Por exemplo: a rocha pode conter menos muscovita; parte da muscovita pode
já ter se decomposto abaixo do solidus, com quartzo, em feldspato potássico e
aluminossilicato, e sua água ter saído como fluido antes do início da fusão; o fundido pode conter mais
água que \qty{6}{\percent}; parte da água pode permanecer na biotita.

\noindent\textbf{Parte II.}

\textbf{7.} $680 - 7.5\times(7.8 - 2) = \qty{636.5}{\celsius}$; a rocha tem seu pico
\qty{43.5}{\celsius} abaixo dele e não pode fundir, mesmo com água livre.

\textbf{8.} $605 + (19.2 - 15.3)/(27 - 15.3)\times65 = \qty{626.7}{\celsius}$. Em
$t = 0$, a base, a \qty{595}{\celsius}, está \qty{32}{\celsius} abaixo; a
\qty{10}{Myr}, a \qty{656}{\celsius}, está acima.

\textbf{9.} A fusão no solidus úmido precisa de água livre, e uma rocha metamórfica aquecida
contém muito pouca: a água livre é consumida pelo primeiro fundido, que então para de crescer
até que as micas se decomponham.

\textbf{10.} Entre 30 e \qty{40}{Myr}: $30 + (800 - 795)/(845 - 795)\times10 =
\qty{31.0}{Myr}$.

\textbf{11.} A base desce abaixo de \qty{12.5}{kbar} entre 70 e \qty{80}{Myr}:
$70 + (13.7 - 12.5)/(13.7 - 12.4)\times10 = \qty{79.2}{Myr}$, a
$914 - 14\times0.923 = \qty{901}{\celsius}$, dentro do intervalo do anfibólio: um anfibolito
na base estaria fundindo.

\textbf{12.} $90 + (11.0 - 10)/(11.0 - 9.6)\times10 = \qty{97.1}{Myr}$, a $867 -
53\times0.714 = \qty{829}{\celsius}$. Isso fica acima do intervalo da muscovita: a muscovita
já se decompôs. Fica dentro dos 800 a \qty{850}{\celsius} da decomposição da biotita
(medidos a pressão mais baixa): a biotita, o outro mineral hidratado do xisto, está no seu
limite.

\noindent\textbf{Parte III.}

\textbf{13.} $n(\ce{Al2O3}) = 15.0/101.96 = \qty{0.14712}{mol}$, $n(\ce{CaO}) =
0.55/56.08 = 0.00981$, $n(\ce{Na2O}) = 3.4/61.98 = 0.05486$, $n(\ce{K2O}) = 4.9/94.20 =
0.05202$.

\textbf{14.} A/CNK $= 0.14712/0.11669 = 1.261$.

\textbf{15.} A/NK $= 0.14712/0.10688 = 1.377$; A/CNK $> 1$: \omterm{def:b2:crustal-melting:asi}{peraluminoso}.

\textbf{16.} \omterm{def:b2:crustal-melting:asi}{Peraluminoso}, vindo de um xisto: um \omterm{def:b2:crustal-melting:types}{granito tipo S}. Seu excesso de alumínio fica na
muscovita e na granada (ou na cordierita).

\textbf{17.} $0.14712 - 0.11669 = \qty{0.03043}{mol}$ de \ce{Al2O3}; na forma de muscovita,
$0.03043\times398.33 = \qty{12.1}{\gram}$ em \qty{100}{\gram} de granito.

\textbf{18.} A obsidiana é (por pouco) \omterm{def:b2:crustal-melting:asi}{metaluminosa}, sem o excesso de alumínio de um
fundido metassedimentar: sua fonte é provavelmente diferente. Mas uma amostra prova pouco:
toda suíte de granitoides tem membros ligeiramente \omterm{def:b2:crustal-melting:asi}{peraluminosos}, e uma suíte se julga por sua
tendência inteira.

\noindent\textbf{Parte IV.}

\textbf{19.} $10^{12}\,\unit{\cubic\metre}\times2650 = \qty{2.65e15}{kg}$.

\textbf{20.} O granito é o fundido menos sua água: $2.65\times10^{15}/0.94 =
\qty{2.82e15}{kg}$ de fundido.

\textbf{21.} Apenas $0.151\times0.70$ da massa do xisto sai como fundido:
$2.82\times10^{15}/(0.151\times0.70) = \qty{2.67e16}{kg}$.

\textbf{22.} $2.67\times10^{16}/2800 = \qty{9.5e12}{\cubic\metre}$, isto é,
\qty{9.5e3}{\cubic\kilo\metre}, quase dez vezes o \omterm{def:b2:magma-emplacement:bodies}{plúton}.

\textbf{23.} $9.5\times10^3/5 = \qty{1.9e3}{\square\kilo\metre}$, um quadrado de cerca de
\qty{44}{km} de lado.

\textbf{24.} Com $F = \qty{7.21}{\percent}$, a fonte é $15.1/7.21 = 2.1$ vezes
maior, cerca de \qty{2.0e4}{\cubic\kilo\metre}. O volume de crosta-fonte fundida para produzir
o \omterm{def:b2:magma-emplacement:bodies}{plúton}, com a \omterm{def:b2:mantle-melting:fraction}{fração de fundido} da questão~5, é de
\textbf{\qty[text-series-to-math=true]{9.5e3}{\cubic\kilo\metre}}.
\end{solution}
